\documentclass[12pt,a4paper]{article}
\usepackage[utf8]{inputenc}
\usepackage[vietnamese]{babel}
\usepackage[margin=2cm]{geometry}
\usepackage{tcolorbox}
\usepackage{enumitem}

\pagestyle{plain}

\begin{document}

% --- HEADER 2 CỘT NGHỊ ĐỊNH 30/2020/NĐ-CP ---
\noindent
\begin{minipage}[t]{0.45\textwidth}
    \begin{center}
        \textbf{CÔNG AN THÀNH PHỐ}\\
        \textbf{\underline{PHÒNG CẢNH SÁT HÌNH SỰ}}\\[0.2cm]
        \small{Số: 07A/BB-LK}\\
        \textit{V/v lấy lời khai đối tượng Trần Ngọc Mai}
    \end{center}
\end{minipage}
\hfill
\begin{minipage}[t]{0.52\textwidth}
    \begin{center}
        \textbf{CỘNG HÒA XÃ HỘI CHỦ NGHĨA VIỆT NAM}\\
        \textbf{\underline{Độc lập - Tự do - Hạnh phúc}}\\[0.2cm]
        \textit{Hà Nội, ngày 25 tháng 07 năm 2026}
    \end{center}
\end{minipage}

\vspace{0.8cm}

\begin{center}
    \textbf{\LARGE BIÊN BẢN GHI LỜI KHAI}\\
    \textbf{\large (ĐỐI TƯỢNG: TRẦN NGỌC MAI - CHUYÊN ÁN TEST-99)}
\end{center}

\vspace{0.3cm}

\begin{tcolorbox}[colback=gray!8!white,colframe=black!75!black,title=\textbf{I. THÔNG TIN NGƯỜI ĐƯỢC LẤY LỜI KHAI}]
\begin{itemize}[leftmargin=*]
    \item \textbf{Họ và tên:} Trần Ngọc Mai | \textbf{Năm sinh:} 1990 (36 tuổi)
    \item \textbf{CCCD số:} 03609000xxxx | \textbf{Nghề nghiệp:} Nhân viên kế toán
    \item \textbf{Mối quan hệ với nạn nhân:} Em họ (Con gái của cô ruột nạn nhân)
    \item \textbf{Cán bộ lấy lời khai:} Thượng úy Trần Quốc Bảo - Điều tra viên
\end{itemize}
\end{tcolorbox}

\section*{II. NỘI DUNG HỎI VÀ ĐÁP CHI TIẾT}

\begin{description}[leftmargin=1.5cm,style=nextline]
    \item[\textbf{Hỏi (Điều tra viên):}] Bà có mối quan hệ như thế nào với nạn nhân Nguyễn Văn Khang?
    \item[\textbf{Đáp (Trần Ngọc Mai):}] Tôi là em họ của anh Khang. Bố mẹ anh Khang mất sớm, hai anh em cùng lớn lên trong khu Bờ Sông. Sau khi ông nội qua đời cách đây 3 tháng, giữa tôi và anh Khang xảy ra tranh chấp về quyền thừa kế căn nhà số 14.

    \item[\textbf{Hỏi (Điều tra viên):}] Chiều tối ngày 24/07/2026, bà có đến căn nhà số 14 Đường Bờ Sông không? Đến lúc mấy giờ và rời đi lúc mấy giờ?
    \item[\textbf{Đáp (Trần Ngọc Mai):}] Khoảng 18:30 chiều 24/07, chồng tôi là anh Lê Quang Vũ lái xe đưa tôi đến nhà anh Khang để nói chuyện rõ ràng về tờ di chúc. Tôi ở đó đến khoảng 19:00 thì rời đi.

    \item[\textbf{Hỏi (Điều tra viên):}] Nội dung cuộc nói chuyện giữa bà và nạn nhân là gì? Tại sao hàng xóm lại nghe thấy tiếng cãi vã lớn?
    \item[\textbf{Đáp (Trần Ngọc Mai):}] Anh Khang đưa ra tờ di chúc viết tay bảo ông nội để lại toàn bộ nhà đất đền bù cho anh ấy. Tôi biết tính anh Khang gian giảo nên vô cùng bức xúc. Tôi khẳng định di chúc đó bị anh ấy làm giả nên hai bên có to tiếng cãi qua cãi lại.

    \item[\textbf{Hỏi (Điều tra viên):}] Bà có động chạm hoặc lấy thứ gì trong khoang tủ âm tường của nạn nhân không?
    \item[\textbf{Đáp (Trần Ngọc Mai):}] Lúc anh Khang đi xuống bếp pha nước, tôi có mở tủ âm tường để tìm tờ di chúc. Tôi thấy tờ di chúc nằm trong hộp sắt nên đã lấy ra xem. Thấy chữ viết có dấu hiệu bất thường, tôi lén tráo bản sao vào và mang bản gốc đi.

    \item[\textbf{Hỏi (Điều tra viên):}] Sau khi rời nhà nạn nhân lúc 19:00, bà đã đi đâu và làm gì? Có ai làm chứng không?
    \item[\textbf{Đáp (Trần Ngọc Mai):}] Chồng tôi đưa tôi thẳng đến Văn phòng Luật sư Nguyễn Văn Minh ở trung tâm thành phố. Tôi làm việc với Luật sư Minh từ 19:30 đến 20:30 để nhờ ông làm thủ tục giám định di chúc. Luật sư Minh có thể làm chứng cho tôi.
\end{description}

\vspace{0.8cm}

\noindent
\begin{minipage}[t]{0.45\textwidth}
    \textbf{\textit{Người khai:}}\\
    \textit{(Ký và ghi rõ họ tên)}\\[1.5cm]
    \textbf{Trần Ngọc Mai}
\end{minipage}
\hfill
\begin{minipage}[t]{0.5\textwidth}
    \begin{center}
        \textbf{CÁN BỘ LẤY LỜI KHAI}\\
        \textit{(Ký và ghi rõ họ tên)}\\[1.5cm]
        \textbf{Thượng úy Trần Quốc Bảo}
    \end{center}
\end{minipage}

\end{document}
